%!TEX root = ../thesis.tex

\thispagestyle{plain}

\section*{Abstract}

% BACHELOR'S THESIS ONLY: AStuPO §21 (7).4: There MUST be at least ONE abstract. It MUST be in GERMAN if the Bachelor's Thesis is not written in German.

% MASTER'S THESIS ONLY: AStuPO §21 (7).4: There MUST be at least ONE abstract. It MUST be in GERMAN OR ENGLISH if the Master's Thesis is not written in German.

Session hijacking, the theft and reuse of web session cookies to impersonate victims, remains a major threat in web security. The PUSH (PUF against Session Hijacking) protocol showed that binding a session to a PUF-based hardware dongle is a viable defense, but its threat model trusts the host machine to mediate between the browser and the dongle, which exposes the dongle to a silent signing attack. Furthermore, it relies on a weak SRAM-PUF whose limited challenge-response space forces key reuse.

We propose REPUSH (Reconfigurable PUF against Session Hijacking), which extends PUSH to cover its limitations. REPUSH assumes the entire host to be compromised and introduces a trusted display on the dongle so that sensitive requests require explicit user approval, thereby eliminating silent signing for sensitive requests. Moreover, the SRAM-PUF is replaced with a logically reconfigurable PUF (LR-PUF) built around a weak delay-based CHOICE-PUF, whose reconfiguration property yields a practically unbounded key space and supports a native key rotation mechanism at the end of every session. 

We implement a proof-of-concept dongle combining an FPGA, which hosts the LR-PUF, with an Arduino Nano board, which handles host communication, cryptographic operations and the trusted display. The used CHOICE-PUF achieves a mean of 51.04\% per-bit uniformity and 99.61\% reliability. The LR-PUF queries with error-correction complete in 24.8 ms, and key enrollment takes 50.5 ms on average. Further latency is introduced by the slow cryptographic operations performed by the Arduino's software implementation of Ed25519, which makes signing take around 2.8 seconds, and enrollment and key rotation take around 12 seconds. Signing with REPUSH is slower than PUSH's 419 ms signing latency. 

Our security analysis shows that REPUSH eliminates the silent signing attack for sensitive requests and resolves the enrollment hijacking attack that previously arose from PUSH's trusted-host assumption. The LR-PUF removes the key-reuse limitation inherent to the SRAM-PUF. None of the added latency is perceptible to the user in practice, since non-sensitive requests continue to be served while a signature is pending and enrollment or rotation occurs infrequently. We conclude that REPUSH could be deployed practically, though some limitations remain, such as the exclusion of physical attacks on the dongle from the threat model. Finally, we discuss how future work could solve these limitations. 


%This thesis introduces the REPUSH (Reconfigurable PUF against Session Hijacking) protocol, a continuous authentication protocol that binds a web session to an external hardware dongle to protect against session hijacking, which is the theft and reuse of session cookies to impersonate some user. REPUSH extends on the PUSH protocol, which was first to show that PUF-based hardware binding is a viable defense against session hijacking. PUSH did however come with some limitations that REPUSH aims to solve. Unlike PUSH, REPUSH's threat model assumes that entire host is compromised and introduces a trusted display on the dongle so that sensitive operations require explicit user approval. PUSH uses a weak SRAM-PUF as the primitive for key generation, which severely limits the number of possible keys and forces keys to eventually be reused. REPUSH solves this by using a Logically Reconfigurable PUS construction built on top of a delay-based CHOICE-PUF, which can be reconfigured to radically change the challenge-response behavior and provide an unbounded number of keys. We first give an introduction to session hijacking and the shortcomings of existing defenses, followed by the necessary background on PUFs, reconfigurable PUFs and the PUSH protocol. We then summarize related work before presenting the design of the RPUSH protocol and its proof-of-concept dongle, which combines an FPGA hosting the LR-PUF and an Arduino Nano baord handling host communication, cryptographic operations and the trusted display. We then describe the implementation of each of the protocol's parties and finally evaluate the security and performance of REPUSH by comparing it to PUSH. Finally, the thesis concludes with answers to the posed research questions and a discussion of directions for future work.
